% !TeX spellcheck = ru_RU
% !TEX root = vkr.tex

\section{Обзор литературы}
\label{sec:relatedworks}

Обзор посвящён оцениванию параметров конечных смесей одномерных непрерывных распределений и выбору их структуры. Материал следует этапам решения задачи: модель, оценивание параметров, начальное приближение, выбор структуры. Затем сопоставляются программные инструменты, описываются исходное состояние \texttt{pysatl-mpest} и прикладные сценарии. В конце приведены анализ и выводы.

\subsection{Методика обзора}
\label{sec:review-method}

Цель обзора --- определить методы и программные решения, необходимые для развития \texttt{pysatl-mpest}, и условия их проверки. Обзор отвечает на четыре вопроса:
\begin{enumerate}
    \item какие методы оценивания применимы при заданной структуре смеси и какие их свойства существенны для библиотеки;
    \item как строить начальные приближения и как выбирать число компонент и семейства;
    \item какие возможности для одномерных, в том числе разнородных, смесей уже реализованы в программных средствах и чего в них не хватает;
    \item от какого исходного состояния развивается библиотека и какие опубликованные примеры пригодны для её проверки.
\end{enumerate}

Источниками служат книги, журнальные и конференционные публикации, препринты, документация и исходный код программных средств. Поиск выполнялся с 9 по 21 октября 2026 года в OpenAlex, CrossRef, arXiv и на сайтах журналов и архивов (Project Euclid, The R Journal, Journal of Statistical Software, HAL, PMLR, JMLR). Он состоял из трёх этапов: общие работы, оценивание, инициализация и выбор структуры (54 запроса); методы машинного обучения для выбора числа компонент и семейств (14 запросов); предсказание структуры смеси по признакам выборки (13 запросов). От классических работ выполнялся переход к связанным алгоритмам и реализациям. Реестр составил 445 работ, из них 329 связаны хотя бы с одним требованием к библиотеке или прикладным сценарием; в тексте использован 41 источник. Ссылки проверены по состоянию на 26 октября 2026 года.

Включаются работы об оценивании конечных смесей одномерных непрерывных распределений и общие методы, имеющие применимый в этой области частный случай. Исключаются многомерные наблюдения, дискретные компоненты, регрессия и временная зависимость. Если доступна только аннотация, источник используется лишь для подтверждаемых ею утверждений.

Для сравнения программных средств составлен каталог из 86 требований: оценивание (24), структура (23), инициализация (6) и остальные категории (33) --- инфраструктура, технические требования, визуализация, представление, производительность, диагностика и специализированные возможности. Из них 34 обязательных, 32 желательных и 20 факультативных. Для каждой из 2322 пар <<инструмент --- требование>> записаны статус (выполнено, частично, не выполнено), обоснование и ссылка на документацию или код. Оценки выставлены автором по документации и коду без запуска программ, поэтому сводные числа характеризуют заявленные возможности, а не качество работы. Обзор имеет целевой характер, и выводы об отсутствии возможностей ограничены рассмотренными интерфейсами.

\subsection{Модель конечной смеси и границы задачи}
\label{sec:mixture-model}

Пусть наблюдения $x_1,\ldots,x_n$ образуют независимую выборку из распределения с плотностью
\begin{equation}
    p(x;\vartheta,\mathcal{S}) = \sum_{j=1}^{K}\pi_j f_{F_j}(x;\theta_j),
    \qquad \pi_j > 0,\qquad \sum_{j=1}^{K}\pi_j=1.
    \label{eq:mixture-density}
\end{equation}
Здесь $K$ --- конечное число компонент, $F_j$ --- параметрическое семейство компоненты, $\theta_j$ --- её параметры, $\pi_j$ --- вес. Структуру модели обозначает $\mathcal{S}=(K,F_1,\ldots,F_K)$, числовые параметры --- $\vartheta=(\pi_1,\ldots,\pi_K,\theta_1,\ldots,\theta_K)$. Каждая компонента задаёт непрерывное распределение одной действительной случайной величины. В однородной смеси семейства совпадают, в разнородной различаются~\cite{McLachlanPeel2000}. При заданной структуре метод максимального правдоподобия максимизирует по допустимым параметрам
\begin{equation}
    \ell(\vartheta;\mathcal{S})=
    \sum_{i=1}^{n}\log\left(\sum_{j=1}^{K}
    \pi_j f_{F_j}(x_i;\theta_j)\right).
    \label{eq:observed-likelihood}
\end{equation}
Выбор структуры --- отдельная задача: возможность задать нормальную и экспоненциальную компоненты не определяет, следует ли предпочесть такую модель двум нормальным.

Идентифицируемость означает восстановимость представления по распределению с точностью до перестановки компонент. Тейхер доказал её для конечных смесей нормальных и гамма-распределений~\cite{Teicher1963}; на произвольное объединение семейств результат не переносится. Экспоненциальное распределение входит в гамма-семейство и семейство Вейбулла как частный случай, поэтому разные наборы семейств могут описывать один закон, и структурный поиск должен учитывать такие эквивалентности. Число компонент также не совпадает с числом мод плотности: две близкие нормальные компоненты образуют одну моду.

Без ограничений задача может быть некорректной: в нормальной смеси масштаб одной компоненты можно устремить к нулю около отдельного наблюдения, и правдоподобие неограниченно возрастает~\cite{RednerWalker1984}. Хэтуэй показал, что ограничение отношений масштабов компонент делает задачу корректной~\cite{Hathaway1985}; используются также нижняя граница масштаба и штраф. Оценка относится к выбранной ограниченной задаче, поэтому при сравнении методов ограничения должны совпадать, а ошибка параметров на синтетических данных вычисляется после сопоставления компонент.

\subsection{Методы оценивания при заданной структуре}
\label{subsec:mixture-estimation-methods}

\subsubsection{Итерационные методы и прямая максимизация}

Скрытая метка компоненты позволяет представить оценивание как задачу с неполными данными. Метод ожидания и максимизации (EM) чередует вычисление условного распределения меток и обновление параметров~\cite{Dempster1977,RednerWalker1984}. На E-шаге при текущих оценках вычисляются апостериорные веса
\[
 r_{ik}^{(t)}=
 \frac{\pi_k^{(t)}f_k(x_i;\theta_k^{(t)})}
 {\sum_{j=1}^{K}\pi_j^{(t)}f_j(x_i;\theta_j^{(t)})},
\]
то есть вероятности принадлежности наблюдения $x_i$ компоненте $k$. По ним строится ожидаемое логарифмическое правдоподобие полных данных
\[
 Q(\boldsymbol{\pi},\boldsymbol{\theta}\mid
     \boldsymbol{\pi}^{(t)},\boldsymbol{\theta}^{(t)})
 =\sum_{i=1}^{n}\sum_{k=1}^{K}r_{ik}^{(t)}
   \bigl[\log\pi_k+\log f_k(x_i;\theta_k)\bigr],
\]
которое на M-шаге максимизируется. Для разнородной смеси меняется только вид $f_k$ у разных компонент.

Точная максимизация $Q$ даёт неубывание наблюдаемого правдоподобия. Обобщённый EM (GEM) сохраняет это свойство, если шаг лишь не уменьшает $Q$~\cite{Dempster1977}, поэтому численный M-шаг должен контролировать именно $Q$, а не внутренний критерий оптимизатора. Монотонность не гарантирует глобального максимума: при условиях регулярности предельные точки стационарны~\cite{Wu1983}, но результат зависит от начальных значений.

Метод ожидания и условной максимизации (ECM) разбивает параметры на блоки и после одного E-шага максимизирует $Q$ по каждому блоку при фиксированных остальных. Последовательность остаётся GEM, а для стационарности блоки должны охватывать всё пространство параметров~\cite{MengRubin1993}. Вариант ECME заменяет часть условных шагов максимизацией наблюдаемого правдоподобия по блоку~\cite{LiuRubin1994}. Монотонность сохраняется лишь при согласованном порядке: шаги по $Q$ выполняются первыми, а после шага по наблюдаемой цели нужен новый E-шаг. Для параметров, меняющих носитель плотности, условия гладкости проверяются отдельно.

Прямая максимизация оптимизирует~\eqref{eq:observed-likelihood} без E-шага~\cite{RednerWalker1984}. Градиентные методы, например L-BFGS-B с границами параметров, требуют производных логарифма правдоподобия; методы Нелдера--Мида и Пауэлла обходятся без них~\cite{ScipyOptimize}. Ограничения учитываются заменой переменных, например экспонентой для масштаба и нормированными экспонентами для весов. Локальные экстремумы и зависимость от старта при этом сохраняются.

\subsubsection{Методы моментов}

Метод моментов приравнивает выборочные моменты теоретическим. Момент смеси линейно выражается через моменты компонент, но после подстановки их выражений система нелинейна относительно весов и параметров. Её решения могут быть неединственными, недопустимыми или чувствительными к ошибкам высоких моментов. У и Ян для одномерной нормальной смеси с общей дисперсией проектируют шумные оценки моментов на допустимое множество; гарантии относятся к этой модели~\cite{WuYang2020}.

L-моменты строятся по порядковым статистикам, существуют при конечном первом абсолютном моменте и менее чувствительны к выбросам~\cite{Hosking1990}. Из их определения следует, что L-моменты смеси зависят от её квантильной функции и, кроме первого, не являются взвешенной суммой L-моментов компонент. Моментное решение не максимизирует ни $Q$, ни наблюдаемое правдоподобие, поэтому в роли M-шага оно требует проверки на увеличение $Q$.

\subsubsection{Аналитические и численные обновления параметров}
\label{subsec:parameter-updates}

Явная формула M-шага существует не для каждого семейства. При различных параметрах компонент $Q$ разделяется по компонентам, а вес обновляется как доля суммарного апостериорного веса. Для нормальной и экспоненциальной компонент обновления замкнуты. Для гамма-компоненты и компоненты Вейбулла уравнение по форме решается численно, после чего масштаб обновляется условно; так устроен \texttt{mixR}~\cite{MixRManual}. Численный шаг принимается, если он не уменьшает $Q$: корень уравнения градиента может оказаться седловой точкой или выйти за допустимую область.

Частичное оценивание имеет два смысла: параметр фиксирован на всём протяжении процедуры либо обновляется отдельным блоком при условной фиксации остальных, как в ECM и ECME. Нулевой эффективный вес компоненты и нулевая плотность наблюдения под всеми компонентами требуют явной обработки: повторной инициализации, удаления компоненты или сообщения о несовместимости данных с моделью.

\subsubsection{Сопоставление методов}
\label{subsec:estimation-methods-table}

Таблица~\ref{tab:estimation-methods} сводит предпосылки и ограничения методов. EM, ECM и ECME задают организацию итераций, а аналитическое или численное обновление --- способ выполнения шага. Метод моментов образует отдельную стратегию без гарантий EM.

\begin{table}[!ht]
\centering
\small
\renewcommand{\arraystretch}{1.25}
\caption{Предпосылки и ограничения методов оценивания при заданной структуре}\label{tab:estimation-methods}
\begin{tabularx}{\textwidth}{@{}>{\raggedright\arraybackslash}p{0.22\textwidth}>{\raggedright\arraybackslash}X>{\raggedright\arraybackslash}X@{}}
\toprule
\textbf{Метод} & \textbf{Предпосылки} & \textbf{Ограничения}\\
\midrule
EM & Вычислимы E-шаг и максимум $Q$ & Локальный максимум, медленная сходимость, нужны ограничения против вырождения\\
\addlinespace
ECM & Вычислимы условные максимумы $Q$; блоки охватывают все свободные параметры & Нужны явный порядок блоков и контроль $Q$\\
\addlinespace
ECME & Согласован порядок шагов по $Q$ и по наблюдаемой цели & При неизвестном носителе монотонность требует отдельной проверки\\
\addlinespace
Прямая максимизация & Вычислимо наблюдаемое правдоподобие, заданы ограничения & Локальные экстремумы, зависимость от параметризации и старта\\
\addlinespace
Метод моментов & Конечны выбранные моменты, система разрешима & Неединственные и недопустимые корни, нет гарантий EM\\
\addlinespace
Метод L-моментов & Конечен первый абсолютный момент & Нет линейного смешивания, решение зависит от семейства\\
\bottomrule
\end{tabularx}
\end{table}

Вне сопоставления остались ускорение сходимости EM (Aitken, SQUAREM), стохастический и классификационный EM, покомпонентный EM, оценивание методом минимального расстояния, робастное оценивание и байесовское сэмплирование, в том числе с неизвестным числом компонент~\cite{RichardsonGreen1997}. Эти методы меняют схему оценивания и в каталоге отнесены к желательным и факультативным.

\subsection{Построение начальных приближений}
\label{sec:initialization}

Начальное приближение задаёт веса, параметры и допустимость первой итерации. Карлис и Ксэкалаки сравнивают случайные старты, приближения по моментам и предварительное разбиение выборки для нормальных смесей~\cite{KarlisXekalaki2003}, Бьернацки, Селё и Говер --- стратегии выбора начальных значений для многомерных нормальных смесей~\cite{Biernacki2003}. Милькович и Грюн для положительных одномерных смесей используют инициализацию по расстояниям, методом $k$-средних и случайную, по сто запусков каждой~\cite{MiljkovicGrun2016}. Выводы этих работ не дают ранжирования инициализаторов для произвольных разнородных смесей, но задают постановку сравнения.

В одномерном случае детерминированное приближение строится по квантильным сегментам упорядоченной выборки или кластеризацией. Оба способа опираются на жёсткое разбиение, поэтому при перекрытии компонент и малом весе одной из них дают смещённый старт. Панич, Клеменц и Нагоде строят старт по предварительной оценке плотности методом REBMIX~\cite{Panic2020}. Повторные запуски со случайных стартов с выбором наибольшего правдоподобия образуют мультистарт; его результат зависит от числа запусков, поэтому методы сравниваются при одинаковом бюджете стартов. Отсюда следует отдельный интерфейс инициализации: его результат передаётся любому оценщику, а время построения старта учитывается отдельно.

\subsection{Определение числа компонент и выбор семейств}
\label{sec:structure-selection}

\subsubsection{Критерии сравнения оценённых моделей}

Добавление компоненты расширяет класс плотностей, поэтому одно правдоподобие не задаёт правила выбора структуры. Критерии Акаике (AIC) и байесовский информационный критерий (BIC) записываются как
\begin{equation}
 \mathrm{AIC}=-2\ell(\widehat\vartheta;\mathcal{S})+2d_{\mathcal{S}},
 \qquad
 \mathrm{BIC}=-2\ell(\widehat\vartheta;\mathcal{S})+d_{\mathcal{S}}\log n,
 \label{eq:information-criteria}
\end{equation}
где $d_{\mathcal{S}}$ --- число свободных параметров~\cite{Akaike1974,Schwarz1978}; фиксация параметров его уменьшает. Состоятельность оценки порядка смеси по BIC доказана при условиях на класс плотностей~\cite{Keribin2000}. Нулевые веса и совпадающие компоненты создают нерегулярные случаи, для которых предлагаются модификации BIC~\cite{NguyenNguyen2026}. Критерий интегрированного полного правдоподобия (ICL) добавляет к BIC штраф за неопределённость принадлежности наблюдений~\cite{Biernacki2000} и при перекрытии компонент может предпочесть другую модель.

Тест отношения правдоподобий сравнивает вложенные модели, но распределение хи-квадрат здесь неприменимо: при исчезновении компоненты её вес находится на границе, а параметры не определены. Маклахлан предложил параметрический бутстреп: выборки генерируются из оценённой нулевой модели, на каждой оцениваются обе модели, и строится эмпирическое распределение статистики~\cite{McLachlan1987}. Тест требует множества полных оцениваний, а сбои оптимизации на бутстреп-выборках искажают его калибровку.

\subsubsection{Поиск кандидатов структуры}

Правило сравнения следует отделять от способа получения кандидатов. Перебор наборов из $M$ семейств с повторениями даёт для фиксированного $K$ число $\binom{M+K-1}{K}$ неупорядоченных наборов, поэтому требует ограничения порядка и бюджета. Фигейреду и Джейн включают отбор компонент в оценивание через критерий минимальной длины сообщения~\cite{FigueiredoJain2002}. Вербек, Влассис и Крёсе жадно добавляют нормальные компоненты~\cite{Verbeek2003}, а Уэда и соавторы сливают и расщепляют компоненты при фиксированном порядке~\cite{Ueda2000}. Вариационный байесовский вывод в \texttt{BayesianGaussianMixture} регуляризует веса при заданном максимальном числе компонент, но не даёт распределения порядка~\cite{SklearnBayesian}. Выбор семейства для разнородной модели ни один из этих подходов не решает.

Обучаемые предикторы предлагают число компонент или семейства по признакам выборки, но их цель обучения не совпадает с порядком смеси. Пиментел и де Карвалью рекомендуют число кластеров по индексам качества разбиений $k$-средних~\cite{Pimentel2020}. Хусси и соавторы обучают нейросеть выбирать одно из девяти семейств для одномерной выборки без смесей~\cite{Khoussi2021}. Архитектура Deep Sets обрабатывает неупорядоченные наборы наблюдений~\cite{Zaheer2017}, но вероятностной модели не задаёт. Такие предикторы пригодны как источник гипотез, которые затем проверяются статистическими критериями. Подходы к выбору структуры сведены в таблице~\ref{tab:structure}.

\begin{table}[!ht]
\centering
\small
\renewcommand{\arraystretch}{1.25}
\caption{Подходы к выбору структуры смеси}\label{tab:structure}
\begin{tabularx}{\textwidth}{@{}>{\raggedright\arraybackslash}p{0.26\textwidth}>{\raggedright\arraybackslash}X>{\raggedright\arraybackslash}X@{}}
\toprule
\textbf{Подход} & \textbf{Задача} & \textbf{Ограничения}\\
\midrule
AIC, BIC & Сравнение моделей по подгонке и числу параметров & Нерегулярные случаи при нулевых весах и совпадающих компонентах\\
\addlinespace
ICL & Сравнение с учётом неопределённости принадлежности & Предпочитает более разделённое разбиение\\
\addlinespace
Бутстреп отношения правдоподобий & Проверка вложенных моделей & Множество полных оцениваний; сбои оптимизации влияют на калибровку\\
\addlinespace
Минимальная длина сообщения & Отбор компонент в ходе оценивания & Для разнородных смесей нужен вывод обновлений\\
\addlinespace
Жадное добавление & Добавление компоненты из кандидатов & Выбор семейства не решён\\
\addlinespace
Слияние и расщепление & Выход из локальных решений & Порядок смеси не меняет\\
\addlinespace
Вариационный байесовский вывод & Регуляризация весов & Нет явного порядка, семейства не выбираются\\
\addlinespace
Предикторы машинного обучения & Гипотезы о порядке и семействах & Цель обучения не совпадает с порядком смеси\\
\bottomrule
\end{tabularx}
\end{table}

\subsection{Разнородные смеси в публикациях}
\label{sec:heterogeneous-literature}

Разнородные смеси описаны реже однородных. Коретто и Хенниг рассматривают оценивание максимального правдоподобия для смеси нормальных и равномерных компонент~\cite{Coretto2011}. Би предлагает ECME-алгоритм для смеси логнормального распределения и распределения Парето с неизвестным нижним порогом последнего и отдельно изучает выбор модели~\cite{Bee2026}. Милькович и Грюн сравнивают шесть семейств в однородных смесях, а разнородную смесь Вейбулла и Burr показывают отдельным примером~\cite{MiljkovicGrun2016}. Эти работы относятся к конкретным сочетаниям семейств; общая процедура оценивания и выбора для произвольного словаря семейств в рассмотренных публикациях не предложена.

\subsection{Существующие программные инструменты}
\label{sec:software}

\subsubsection{Сопоставление по каталогу требований}

Рассмотрено 27 инструментов для оценивания смесей и представления распределений на Python, R, Julia и C++. Многие из них рассчитаны также на многомерные, регрессионные или латентно-классовые модели, поэтому сравнение ограничено требованиями к одномерным непрерывным смесям. Таблица~\ref{tab:tools-matrix} показывает число полностью выполненных требований по категориям и общее число частично выполненных.

\begin{table}[p]
\centering
\small
\setlength{\tabcolsep}{2.5pt}
\renewcommand{\arraystretch}{1.12}
\newcommand{\hd}[2]{\begin{tabular}[b]{@{}c@{}}\textbf{#1}\\\textbf{#2}\end{tabular}}
\caption{Число требований каталога, выполненных инструментами}\label{tab:tools-matrix}
\begin{tabularx}{\textwidth}{@{}>{\raggedright\arraybackslash}Xl*{6}{c}@{}}
\toprule
& & \multicolumn{5}{c}{\textbf{Выполнено полностью}} & \\
\cmidrule(lr){3-7}
\textbf{Средство} & \textbf{Среда} & \hd{Оцени-}{вание} & \hd{Струк-}{тура} & \hd{Инициа-}{лизация} & \textbf{Прочие} & \textbf{Всего} & \hd{Частич-}{но}\\
& & из 24 & из 23 & из 6 & из 33 & из 86 & \\
\midrule
\texttt{brms} & R/C++ & 8 & 10 & 3 & 20 & 41 & 11\\
\rowcolor{gray!12}\texttt{flexmix} & R & 6 & 9 & 3 & 14 & 32 & 14\\
\texttt{pysatl-mpest} & Python & 7 & 9 & 1 & 14 & 31 & 6\\
\rowcolor{gray!12}\texttt{mixR} & R & 6 & 5 & 6 & 13 & 30 & 10\\
\texttt{StepMix} & Python & 6 & 6 & 3 & 13 & 28 & 11\\
\rowcolor{gray!12}\texttt{Rmixmod} & R/C++ & 4 & 5 & 3 & 16 & 28 & 10\\
\texttt{mclust} & R & 4 & 6 & 3 & 14 & 27 & 17\\
\rowcolor{gray!12}\texttt{rebmix} & R/C++ & 5 & 9 & 0 & 13 & 27 & 9\\
\texttt{bnpy} & Python & 6 & 6 & 4 & 10 & 26 & 18\\
\rowcolor{gray!12}\texttt{mixtools} & R & 7 & 5 & 3 & 9 & 24 & 18\\
\texttt{Distributions.jl} & Julia & 3 & 7 & 1 & 13 & 24 & 12\\
\rowcolor{gray!12}\texttt{pomegranate} & Python & 5 & 6 & 3 & 9 & 23 & 16\\
\texttt{mixdist} & R & 7 & 5 & 2 & 9 & 23 & 10\\
\rowcolor{gray!12}\texttt{scikit-learn} & Python & 4 & 4 & 3 & 12 & 23 & 10\\
\texttt{sn} & R & 6 & 4 & 2 & 11 & 23 & 8\\
\rowcolor{gray!12}\texttt{EMCluster} & R/C & 2 & 5 & 3 & 12 & 22 & 9\\
\texttt{CAMAN} & R/C & 4 & 7 & 2 & 7 & 20 & 9\\
\rowcolor{gray!12}\texttt{ClusterR} & R/C++ & 3 & 3 & 2 & 12 & 20 & 4\\
\texttt{distfit} & Python & 2 & 6 & 0 & 11 & 19 & 18\\
\rowcolor{gray!12}\texttt{teigen} & R & 2 & 6 & 2 & 9 & 19 & 14\\
\texttt{GaussianMixtures.jl} & Julia & 2 & 3 & 3 & 10 & 18 & 9\\
\rowcolor{gray!12}\texttt{mlpack} & C++ & 2 & 2 & 1 & 13 & 18 & 7\\
\texttt{Mixture-Models} & Python & 4 & 5 & 3 & 5 & 17 & 5\\
\rowcolor{gray!12}\texttt{fpc} & R & 2 & 5 & 2 & 7 & 16 & 13\\
\texttt{student-mixture} & Python & 4 & 5 & 3 & 3 & 15 & 10\\
\rowcolor{gray!12}\texttt{Armadillo} & C++ & 2 & 2 & 2 & 9 & 15 & 6\\
\texttt{mixem} & Python & 2 & 5 & 1 & 6 & 14 & 5\\
\bottomrule
\end{tabularx}
\end{table}

Суммарные числа не ранжируют качество. Лидирует \texttt{brms} --- байесовский пакет для регрессионных моделей, в котором смесь является частным случаем, а около половины выполненных требований относятся к категориям вне оценивания. В категориях, существенных для работы, видно следующее:
\begin{itemize}
    \item все требования по инициализации полностью выполнены только в \texttt{mixR}, в \texttt{pysatl-mpest} --- одно из шести;
    \item критерии AIC и BIC полностью реализованы в 15 инструментах, параметрический бутстреп отношения правдоподобий --- в семи (\texttt{CAMAN}, \texttt{fpc}, \texttt{mclust}, \texttt{mixR}, \texttt{mixtools}, \texttt{rebmix}, \texttt{StepMix});
    \item замкнутые формулы M-шага ни один инструмент не поддерживает полностью: они существуют для отдельных семейств, остальные случаи решаются численно;
    \item ECME реализован только в \texttt{mixtools} и \texttt{pysatl-mpest}, стратегия по наблюдаемому правдоподобию --- только в \texttt{pysatl-mpest}.
\end{itemize}

\subsubsection{Устройство представительных средств}

Устройство инструментов определяет, как в них добавляется новое семейство. В \texttt{scikit-learn} класс \texttt{GaussianMixture} поддерживает только нормальные компоненты, несколько запусков, регуляризацию дисперсии и критерии AIC и BIC~\cite{SklearnMixture}. В \texttt{mixtools} каждое семейство обслуживается отдельной функцией, а бутстреп-тест числа компонент доступен для перечисленных в интерфейсе моделей~\cite{Benaglia2009,MixtoolsManual}. В \texttt{mixR} аргумент \texttt{family} выбирает одно из четырёх семейств для всей смеси; обновления аналитические для нормального и логнормального семейств и частично численные для гамма и Вейбулла~\cite{Yu2022,MixRManual}. В \texttt{flexmix} итерационный механизм отделён от драйвера компонентной модели~\cite{Leisch2004,GrunLeisch2008}; драйвер \texttt{FLXMCdist1} обслуживает семь положительных семейств, но принимает одно семейство на модель~\cite{FlexmixManual}. Пакет \texttt{rebmix} выбирает число компонент по предварительной оценке плотности и также назначает одно семейство на переменную~\cite{Nagode2015,RebmixManual}. Разнородные смеси строятся в \texttt{pomegranate} из объектов распределений, а в \texttt{mixem} --- из классов с логарифмом плотности и взвешенным оцениванием~\cite{PomegranateMixture,MixemAPI}.

\subsubsection{Ближайшие аналоги по разнородности}

Разнородную смесь допускают пять средств каталога: \texttt{pysatl-mpest}, \texttt{pomegranate}, \texttt{mixem}, \texttt{Distributions.jl} и \texttt{brms}. Их возможности сопоставлены в таблице~\ref{tab:hetero}.

\begin{table}[!ht]
\centering
\small
\setlength{\tabcolsep}{4pt}
\renewcommand{\arraystretch}{1.2}
\caption{Возможности средств, допускающих разнородные смеси ($\checkmark$~--- выполнено, прочерк --- не выполнено)}\label{tab:hetero}
\begin{tabularx}{\textwidth}{@{}>{\raggedright\arraybackslash}X*{5}{>{\centering\arraybackslash}p{0.11\textwidth}}@{}}
\toprule
\textbf{Возможность} & \footnotesize\texttt{pysatl-}\newline\texttt{mpest} & \footnotesize\texttt{pome-}\newline\texttt{granate} & \footnotesize\texttt{mixem} & \footnotesize\texttt{Distribu-}\newline\texttt{tions.jl} & \footnotesize\texttt{brms}\\
\midrule
Блочная схема ECM & $\checkmark$ & --- & --- & --- & ---\\
\rowcolor{gray!12}Схема ECME & $\checkmark$ & --- & --- & --- & ---\\
Фиксация подмножеств параметров & $\checkmark$ & $\checkmark$ & --- & --- & $\checkmark$\\
\rowcolor{gray!12}Явные границы параметров & --- & частично & --- & --- & $\checkmark$\\
Мультистарт & --- & --- & --- & --- & $\checkmark$\\
\rowcolor{gray!12}Кластерная инициализация & --- & $\checkmark$ & --- & --- & ---\\
Критерии AIC и BIC & частично & --- & --- & частично & $\checkmark$\\
\rowcolor{gray!12}Бутстреп отношения правдоподобий & --- & --- & --- & --- & ---\\
График декомпозиции плотности & --- & --- & --- & частично & $\checkmark$\\
\bottomrule
\end{tabularx}
\end{table}

Мультистарт, информационные критерии, фиксация параметров и явные границы одновременно есть только в \texttt{brms}, который использует байесовское сэмплирование, а не итерационную оптимизацию. Бутстреп отношения правдоподобий не реализован ни в одном из пяти средств. ECM и ECME в сочетании с фиксацией подмножеств параметров есть только в \texttt{pysatl-mpest}. Итерационного инструмента для разнородных смесей с мультистартом и выбором структуры среди рассмотренных нет.

\subsection{Технологическая основа и исходное состояние pysatl-mpest}
\label{sec:mpest-state}

Библиотека написана на Python (версии 3.12 и 3.13) с NumPy и SciPy. Качество кода поддерживают \texttt{pytest}, \texttt{hypothesis}, \texttt{ruff}, \texttt{mypy} и \texttt{pre-commit}. Вероятностная часть переводится на ядро \texttt{pysatl-core}, которое предоставляет параметрические семейства с несколькими параметризациями, проверку допустимости параметров, градиент логарифма правдоподобия, частичные семейства с фиксированными параметрами и конечную смесь~\cite{CoreRepository}. Семейство, добавленное в ядро, становится доступно оценщикам без дублирования плотности.

В \texttt{pysatl-mpest} реализованы представление смеси, интерфейс оценщика и итерационный конвейер: E-шаг, блочное обновление параметров, правила остановки и удаление компонент с малым весом~\cite{MpestRepository}. Есть фасады ECM и ECME и стратегии M-шага по $Q$, по наблюдаемому правдоподобию и по методу моментов; универсальная численная реализация дополнена аналитическими обновлениями для нормального и экспоненциального семейств. Используются только безградиентные оптимизаторы Нелдера--Мида и Пауэлла. Градиенты ядра не задействованы, а для градиентных методов с границами нужно сопоставить допустимые области ядра с числовыми интервалами. Модуль прямых оценщиков пуст, прежние модули инициализации удалены при переходе на ядро, выбора структуры нет, часть зависимостей (\texttt{xgboost}, \texttt{kneed}, \texttt{pywavelets}, \texttt{scikit-fuzzy}) не используется.

Из 34 обязательных требований каталога полностью выполнены 26, частично --- три (общий каркас EM, замкнутые формулы M-шага, критерии AIC и BIC). Не выполнены пять: явные границы параметров, мультистарт, кластерная инициализация, параметрический бутстреп и график декомпозиции плотности. Инфраструктура выполнена полностью, а из девяти требований к диагностике и визуализации не выполнено ни одно, хотя график декомпозиции плотности есть у 15 инструментов, а стандартные ошибки и доверительные интервалы --- у 10.

\subsection{Прикладные сценарии}
\label{sec:review-scenarios}

Прикладные сценарии различаются свойствами наблюдений, от которых зависят требования к оцениванию: формой компонент, носителем, тяжестью хвоста, долей малых групп и способом регистрации данных. Десять сценариев каталога приведены в таблице~\ref{tab:scenarios}. Часть из них выходит за модель полностью наблюдаемой независимой выборки: цензурирование меняет правдоподобие, а сгруппированные данные требуют оценивания по интервальным частотам.

\begin{table}[!ht]
\centering
\small
\renewcommand{\arraystretch}{1.2}
\caption{Прикладные сценарии каталога}\label{tab:scenarios}
\begin{tabularx}{\textwidth}{@{}>{\raggedright\arraybackslash}p{0.27\textwidth}l>{\raggedright\arraybackslash}p{0.21\textwidth}>{\raggedright\arraybackslash}X@{}}
\toprule
\textbf{Сценарий} & \textbf{Смесь} & \textbf{Компоненты} & \textbf{Особенности данных}\\
\midrule
Лучевые скорости галактик & Однородная & Симметричные & Малая выборка, число групп неизвестно\\
\rowcolor{gray!12}Кислотность горных озёр & Однородная & Симметричные & Сильное перекрытие, число групп неизвестно\\
Толщина бумаги гербовых марок & Однородная & Симметричные & Перекрытие, число групп неизвестно\\
\rowcolor{gray!12}Интервалы между извержениями гейзеров & Однородная & Симметричные & Разделённые компоненты\\
Активность ферментов & Разнородная & Асимметричные & Перекрытие компонент\\
\rowcolor{gray!12}Проточная цитометрия & Однородная & Асимметричные & Очень малые группы, большая выборка\\
Время до отказа технических систем & Разнородная & Положительный носитель & Цензурирование\\
\rowcolor{gray!12}Скорость ветра & Разнородная & Положительный носитель & Перекрытие, большая выборка\\
Страховые убытки & Разнородная & Тяжёлый хвост & Малый вес хвостовой компоненты\\
\rowcolor{gray!12}Задержки пакетов в сетях & Разнородная & Тяжёлый хвост & Сгруппированные данные, очень большая выборка\\
\bottomrule
\end{tabularx}
\end{table}

Для части сценариев опубликованы примеры с открытыми данными или кодом. Статья о \texttt{mixR} содержит код генерации двухкомпонентной смеси Вейбулла из 1000 наблюдений с известными параметрами~\cite{Yu2022,MixRManual}. Би исследует смесь логнормального распределения и распределения Парето при объёмах выборки 100 и 500, весе логнормальной компоненты от $0{,}7$ до $0{,}95$ и нескольких значениях параметра формы и порога Парето; код доступен в пакете \texttt{LNPar}~\cite{Bee2026,LNPar}. При единичном параметре формы у распределения Парето нет первого момента, и методы моментов неприменимы. Милькович и Грюн подгоняют однородные смеси шести семейств с числом компонент от одного до восьми к 2492 датским страховым убыткам за 1980--1990 годы и приводят значения правдоподобия, AIC и BIC; данные доступны в \texttt{SMPracticals}~\cite{MiljkovicGrun2016,SMPracticalsManual}. На синтетических примерах истинные параметры известны, на реальных данных оцениваются только качество подгонки и устойчивость к старту.

\subsection{Анализ обзора и выводы для развития библиотеки}
\label{sec:review-conclusions}

\subsubsection{Основные результаты обзора}

\begin{enumerate}
    \item EM, ECM, ECME, прямая максимизация и методы моментов различаются целевой функцией, гарантиями и предпосылками. Монотонность EM зависит от контроля $Q$ на каждом шаге, а аналитические обновления существуют лишь для отдельных семейств.
    \item Выбор числа компонент и выбор семейств --- отдельные задачи. Обоснованы информационные критерии и параметрический бутстреп для вложенных моделей; обучаемые предикторы гарантий не дают.
    \item Разнородные смеси допускают пять из 27 инструментов. Итерационного средства с мультистартом, критериями выбора и бутстрепом среди них нет; сочетание разнородных компонент, ECM, ECME и фиксации подмножеств параметров есть только в \texttt{pysatl-mpest}.
    \item Из 34 обязательных требований \texttt{pysatl-mpest} полностью выполняет 26. Пробелы сосредоточены в инициализации, выборе структуры, диагностике и визуализации.
    \item В литературе есть примеры с открытыми данными или кодом для положительной асимметрии, разнородности и тяжёлого хвоста, пригодные для экспериментальной проверки.
\end{enumerate}

\subsubsection{Архитектурные решения и приоритеты}

Граница ответственности следует из разделения модели и алгоритма. \texttt{pysatl-core} отвечает за семейства, параметризации, вероятностные характеристики, градиенты и генерацию данных, а \texttt{pysatl-mpest} --- за состояние оценивания, веса, апостериорные вероятности, стратегии обновления, остановку, инициализацию и поиск структуры.

Целевая функция и способ её оптимизации задаются раздельно (таблица~\ref{tab:estimation-methods}). Стратегия получает семейство, свободные параметры, ограничения и состояние E-шага; аналитическое обновление применяется при выполнении его предпосылок, иначе выполняется численная условная оптимизация. Приём шага проверяется по соответствующей цели, порядок блоков ECM и ECME задаётся явно. Градиенты ядра вместе с явными границами параметров позволяют подключить градиентные оптимизаторы.

Инициализация выделяется в отдельный модуль, который возвращает допустимую модель и состояние генератора; базовый набор включает пользовательский старт, квантильное разбиение, кластеризацию и мультистарт. Для выбора структуры основным правилом служит BIC при ограниченном переборе числа компонент и словаря семейств. AIC и ICL используются для сопоставления, параметрический бутстреп --- для вложенных пар моделей с тем же оценщиком и бюджетом стартов. Жадные переходы и удаление компонент сокращают перебор, обучаемые предикторы служат источником кандидатов. Для выпуска библиотеки нужны средства диагностики (график декомпозиции плотности, стандартные ошибки, проверки согласия) и очистка неиспользуемых зависимостей; инфраструктура сборки и тестирования уже соответствует каталогу.

\subsubsection{Соответствие постановке задачи}

Цель работы --- развитие \texttt{pysatl-mpest} для оценивания параметров и выбора структуры смесей одномерных непрерывных распределений с публикацией релиза --- обзором подтверждается. Задачи об интеграции \texttt{pysatl-core} с рефакторингом архитектуры, о модуле выбора структуры, о сборе требований и реализации недостающей функциональности и о проведении экспериментов опираются на описанные выводы. Обзор уточняет содержание задачи о недостающей функциональности: в неё входят явные границы параметров и мультистарт, не выполненные по каталогу обязательные требования, а также средства диагностики, в первую очередь график декомпозиции плотности. Эти средства входят и в публикуемый релиз. Изменения набора задач обзор не требует.

\subsubsection{Принципы сравнения методов}

Корректное сравнение методов требует одинаковых стартов, ограничений и бюджета запусков, а время построения старта учитывается отдельно. Ошибка параметров вычисляется после сопоставления компонент, неверное число компонент считается отдельным исходом. Для тяжёлого хвоста метрика выбирается с учётом существования моментов. Аналитические и численные обновления сравниваются для одной цели и параметризации, причём число итераций не характеризует стоимость: шаг оптимизатора содержит много вычислений плотности.

\subsubsection{Ограничения выводов и ожидаемый вклад}

Выводы о возможностях инструментов основаны на документации и коде без запуска программ, а оценки выставлены одним автором. Опубликованные результаты переносятся только на исследованные в них модели.

По обзору нельзя заявлять новизну EM, ECME, разнородных смесей или общего интерфейса расширения: такие методы и решения уже существуют. Ожидаемый вклад --- согласованное развитие \texttt{pysatl-mpest} на основе \texttt{pysatl-core}: сочетание сменных стратегий, управления свободными параметрами, независимой инициализации и проверяемого поиска структуры с воспроизводимым сравнением обновлений. Это инженерный вклад; вывод о научной новизне отдельного метода потребовал бы его самостоятельного определения и сравнения с ближайшими работами.
